\chapter{A Tour of the Five Libraries}
\label{ch:fivelibs}

\chapterorient{The previous chapter drew the architecture of the synthesized
library---five modules, a topological order, an \texttt{\#extern} boundary where
an opaque kernel sits. This chapter walks the shelves. We open four of the five
libraries one at a time---the shared \emph{types}, and the three calculus
libraries \emph{iteration}, \emph{data-algebra}, and \emph{coordination}---and
read the actual synthesized code each holds: the record definitions, the def
bodies, the \texttt{where}-local leaves, the kernel callouts. (The fifth library,
\emph{drivers}, composes these four into the five analyses and is the subject of
\cref{ch:composing}.) Two def bodies get a line-by-line reading. By the end you
will have seen the whole vocabulary of the simulator rendered as code, and you
will have noticed the one fact the tour exists to make vivid: across all four
libraries the same handful of shapes keeps reappearing, because the library is
small for the same reason the calculus is small.}

\section{Four libraries, one tour}

\Cref{ch:synthlib} laid out the synthesized library as five modules in a
topological stack: a foundational \emph{types} module of shared records, the
three calculus modules---\emph{iteration}, \emph{data-algebra},
\emph{coordination}---and a top \emph{drivers} module that composes them into the
five simulation analyses. A module appears in the stack \emph{after} everything it
uses: \emph{types} defines records the others thread; \emph{iteration} holds the
loop combinator the solvers fold; \emph{data-algebra} holds the tensor verbs the
steps call; \emph{coordination} holds the outer caps that drive the loops; and
\emph{drivers} sits on top, wiring the rest into electrostatics, the driven sweep,
the eigenmode solve, and so on.

This chapter tours the four lower libraries---the shared types and the three
calculus libraries---and shows what each \emph{holds}, in real synthesized code.
\Cref{fig:shelves} is the map for the tour: four labelled shelves, each stocked
with the operators and types we will read. Keep it in view; every section that
follows pulls items off one of these shelves and renders the def body. We save
\emph{drivers} for \cref{ch:composing}, because a driver is not a new kind of
thing---it is a \emph{composition} of the things on these four shelves, and you
cannot compose a vocabulary you have not yet met.

\begin{figure}[t]
\centering
\begin{tikzpicture}[font=\footnotesize, node distance=4mm]
  \def\shelfw{34mm}
  % ---- types ----
  \node[stage, minimum width=\shelfw, align=center, fill=simBgGray, draw=simGray]
    (tH) at (0,0) {\textbf{types}\\\scriptsize shared records};
  \node[opbox, minimum width=\shelfw, align=left, below=of tH] (tB)
    {\code{IoData} \scriptsize\itshape (config)\\
     \code{OpParams} \scriptsize\itshape (build-time)\\
     \code{SimState} \scriptsize\itshape (run-time)};
  % ---- iteration ----
  \node[stage, minimum width=\shelfw, align=center, fill=simBgGreen, draw=simGreen,
    right=9mm of tH] (iH) {\textbf{iteration}\\\scriptsize loops \& steps};
  \node[opbox, minimum width=\shelfw, align=left, below=of iH] (iB)
    {\code{iterate\_while}\\
     \code{iterate\_while\_with\_prev}\\
     \code{krylov\_step} \scriptsize\itshape (A/B)\\
     \code{chebyshev}\\
     \scriptsize\itshape \code{Krylov} \code{StepOutputs} \code{PrevCarry}};
  % ---- data-algebra ----
  \node[stage, minimum width=\shelfw, align=center, fill=simBgGold, draw=simGold,
    right=9mm of iH] (dH) {\textbf{data-algebra}\\\scriptsize tensor verbs};
  \node[opbox, minimum width=\shelfw, align=left, below=of dH] (dB)
    {\code{linear\_combination}\\
     \code{inner\_product} \code{dot} \code{nrm2}\\
     \code{mk\_matrix\_free\_operator}\\
     \code{fe\_assemble} \scriptsize\itshape (\#extern)\\
     \scriptsize\itshape gram / energy / S-param / \dots};
  % ---- coordination ----
  \node[stage, minimum width=\shelfw, align=center, fill=simBgBlue, draw=simBlue,
    right=9mm of dH] (cH) {\textbf{coordination}\\\scriptsize outer caps};
  \node[opbox, minimum width=\shelfw, align=left, below=of cH] (cB)
    {\code{Solve} \scriptsize\itshape monad\\
     \code{ksp\_solve} \code{eigsolve} \scriptsize\itshape(\#extern)\\
     \code{solve\_family}\\
     \code{frequency\_sweep} \code{fold\_solve}\\
     \scriptsize\itshape \code{Outcome} \code{EigState}};
\end{tikzpicture}
\caption[The four libraries side by side]{The four libraries this chapter tours,
each a shelf stocked with the operators and types we will render. \emph{types}
holds the three cross-cutting records every solve threads. \emph{iteration} holds
the loop combinator of \cref{ch:fold} and the step kernels it folds.
\emph{data-algebra} holds the tensor-combining verbs of \cref{ch:combinators}.
\emph{coordination} holds the \code{Solve}-monadic outer caps that drive the
loops over families. The fifth library, \emph{drivers}, composes all four into the
five analyses and is \cref{ch:composing}. Read left to right: the order is
topological---each shelf uses only the shelves to its left.}
\label{fig:shelves}
\end{figure}

One convention governs every code block below, and it is worth stating once. These
libraries are the \emph{implementation view} of the spec: they render the
synthesized code, but they do not restate the laws or the semantics, which live in
the calculus chapters. So each def we read here links \emph{up} to its authoritative
home---\code{iterate\_while}'s laws are in \cref{ch:fold}, \code{linear\_combination}'s
in \cref{ch:combinators}, and so on. The code is what is new; the meaning we have
already built. With that, we open the first shelf.

\section{The \texttt{types} library: three records, three lifetimes}
\label{sec:types}

The foundational shelf holds the records that \emph{every} other library threads.
There are exactly three that are genuinely cross-cutting---shared across two or
more of the calculus libraries---and they are the three lifetimes we named in
\cref{ch:strata}: a configuration parsed once and read forever, a bundle of
operator parameters fixed at construction, and a solve-state that evolves as the
loop runs. \Cref{fig:typecards} draws them as three cards, each tagged with its
stratum.

\begin{figure}[t]
\centering
\begin{tikzpicture}[font=\footnotesize, node distance=6mm]
  \def\cardw{40mm}
  % IoData card
  \node[draw=simGray, fill=simBgGray, rounded corners=3pt, thick,
    minimum width=\cardw, align=left, inner sep=5pt] (io) at (0,0)
    {\textbf{\code{IoData}}\\[1pt]\scriptsize\itshape parse-time, read-only\\[3pt]
     \scriptsize problem / model / domains\\
     \scriptsize boundaries / solver / units\\[2pt]
     \scriptsize\textbf{selects the driver}\\\scriptsize via \code{problem.type}};
  \node[band, fill=simGray, text=white, font=\scriptsize\bfseries, above=1mm of io]
    {stratum 0: parse once};
  % OpParams card
  \node[draw=simGreen, fill=simBgGreen, rounded corners=3pt, thick,
    minimum width=\cardw, align=left, inner sep=5pt, right=10mm of io] (op)
    {\textbf{\code{OpParams}}\\[1pt]\scriptsize\itshape build-time, read-only\\[3pt]
     \scriptsize \code{T} (apply surface)\\
     \scriptsize \code{eps} (stopping test)\\
     \scriptsize variant selectors,\\\scriptsize tolerances, \code{max\_it}};
  \node[band, fill=simGreen, text=white, font=\scriptsize\bfseries, above=1mm of op]
    {stratum 1: fixed at build};
  % SimState card
  \node[draw=simRust, fill=simBgRust, rounded corners=3pt, thick,
    minimum width=\cardw, align=left, inner sep=5pt, right=10mm of op] (ss)
    {\textbf{\code{SimState}}\\[1pt]\scriptsize\itshape run-time, evolving\\[3pt]
     \scriptsize \code{x} (the iterate)\\
     \scriptsize \code{it} \code{converged}\\
     \scriptsize \code{final\_res} \code{initial\_res}\\[2pt]
     \scriptsize\textbf{uniform across solvers}};
  \node[band, fill=simRust, text=white, font=\scriptsize\bfseries, above=1mm of ss]
    {stratum 2: evolves each step};
\end{tikzpicture}
\caption[The \texttt{types} records as stratum-tagged cards]{The three
cross-cutting records of the \emph{types} library, each tagged with its stratum
from \cref{ch:strata}. \code{IoData} is parsed once from the user's config and
read for the rest of the run; its \code{problem.type} field selects which driver
fires. \code{OpParams} is fixed at solver construction and read but never written
during the solve. \code{SimState} is the one record that \emph{evolves}---a fresh
value each step---and it has the \emph{same five fields} for every Krylov-shaped
solver in the book. Three records, three lifetimes; the type-placement rule keeps
only these cross-cutting three here.}
\label{fig:typecards}
\end{figure}

\subsection{\texttt{IoData}: the configuration parsed once}

The first record is the whole configuration, parsed once at startup from the
user's JSON file and then read---never written---for the rest of the run. It is an
aggregate of six sub-records: the problem type and solver-pipeline knobs, the mesh
and material model, the per-domain data, the boundary conditions, the solver
settings, and a units converter.

\begin{lstlisting}[style=l4style]
-- The aggregate parsed once at startup; read-only across the whole solve.
IoData = {
  problem    : ProblemConfig,   -- driver selector (problem.type) + pipeline knobs
  model      : ModelConfig,     -- mesh file + refinement + material assignment
  domains    : DomainConfig,    -- per-domain materials + postprocess energy regions
  boundaries : BoundaryConfig,  -- BC surfaces (PEC/PMC/impedance/lumped-/wave-port)
  solver     : SolverConfig,    -- linear/eigen/driven/transient settings + tolerances
  units      : Units            -- SI <-> nondimensional scale converter
}
-- utility API (parse-time only)
parseConfig :: FilePath -> IoData        -- parse the JSON config tree once
problemType :: IoData -> ProblemType     -- the driver-dispatch selector (a projection)
\end{lstlisting}
The load-bearing field is \code{problem.type}. It is read once, at the very top of
the run, by the dispatch that chooses which of the five drivers to invoke---the
single point where the ``six analyses, one machine'' of \cref{ch:targets} becomes
a concrete branch. The function \code{problemType} that reads it is a trivial
projection: \code{IoData} carries no behaviour, only data. Everything an analysis
does with the configuration, it does by reading fields out of this one immutable
object.

\subsection{\texttt{OpParams}: the operator's build-time parameters}

The second record holds the parameters fixed when the operator is
\emph{constructed}---the variant selectors, the tolerances, and the
constructed-operator surfaces the kernel will apply. It is read by the step kernel
(through its surfaces) and captured once by the outer caps; it never changes during
the solve.

\begin{lstlisting}[style=l4style]
-- read-only; captured once at solve construction, never re-inspected per step.
OpParams = {
  T          : ConstructedOp,    -- the apply surface (preconditioned operator)
  orthog?    : OrthogSurface,    -- GMRES/Arnoldi only; absent (no-op) for CG
  scalars?   : ScalarSurface,    -- Chebyshev polynomial-recurrence scalars; else absent
  eps        : Convergence,      -- the stopping-predicate surface
  -- variant selectors (closed over by the surfaces above; not read by the body)
  pc_side    : PreconditionerSide,
  gs_orthog  : Orthogonalization,
  flexible   : Bool,
  poly_kind? : PolynomialKind,
  restart    : RestartMode,
  -- termination knobs (close into eps)
  max_dim : Int, max_it : Int, rel_tol : Scalar, abs_tol : Scalar
}
\end{lstlisting}
Notice what \cref{ch:strata} promised: the variant selectors (\code{pc\_side},
\code{gs\_orthog}, \code{flexible}, \dots) are \emph{folded into the surfaces}
\code{T}, \code{orthog}, \code{eps} at construction. The kernel never branches on
``am I CG or GMRES?''---that question is answered once, when the surfaces are
built, and the body just applies \code{T} and \code{eps} as given. The
question-marks (\code{orthog?}, \code{scalars?}, \code{poly\_kind?}) mark fields
present for some solvers and absent for others: CG has no orthogonalization
surface, Chebyshev has no Krylov restart. A solver \emph{is} a choice of which of
these are present.

\subsection{\texttt{SimState}: the run-time solve state}

The third record is the only one that evolves. It carries the externally-visible
quantities a Krylov-shaped solve produces and reports, and---this is the fact worth
underlining---it has the \emph{same five fields for every solver in the book}.

\begin{lstlisting}[style=l4style]
-- the value threaded by `Solve a = StateT SimState Identity a`; every field run-time.
-- the iterate `x` is named with shape group S, not a fixed rank-1 axis.
SimState = {
  x           : Tensor[(S: ...)],  -- the current iterate (the solve's primary product)
  it          : Int,               -- iteration count
  converged   : Bool,              -- convergence flag
  final_res   : Scalar,            -- final (absolute) residual, possibly an estimate
  initial_res : Scalar             -- initial (absolute) residual, captured at entry
}
\end{lstlisting}

\begin{keyidea}
CG, GMRES, FGMRES, and Chebyshev all evolve the \emph{same} \code{SimState}: an
iterate \code{x}, a counter \code{it}, a converged flag, and two residuals. The
solvers differ in their \emph{ephemeral} workspace---CG's single search direction,
GMRES's growing basis---but the run-time state they expose to the outside world is
one five-field record. This is why \code{ksp\_solve} can be \emph{one} cap over
\emph{four} solvers (\cref{sec:coordination}): the state it threads is uniform, and
the differences live below it, in the step kernel's private \code{Krylov} bundle.
\end{keyidea}

\begin{notationbox}
Only these three records live in \emph{types}, because only these three are shared
across two or more libraries. A record used by just \emph{one} library---CG's
\code{Krylov} workspace, the eigensolver's \code{EigState}---is placed with that
library instead, bundled next to the operators that consume it. The
\emph{type-placement rule} (\cref{ch:synthlib}) is exactly this: cross-cutting
records go in \emph{types}; single-library records go with their library. We meet
the single-library records as we reach each shelf.
\end{notationbox}

\section{The \texttt{iteration} library: the loop and its kernels}
\label{sec:iteration}

The second shelf holds the heart of the framework: the one loop combinator of
\cref{ch:fold} and the step kernels it folds. Four operators in topological
order---\code{iterate\_while} first (everything below it folds it), then its
two-term-recurrence sibling, then the Krylov step, then the Chebyshev smoother---and
three single-library workspace types rendered just ahead of them. This is the
library whose first def we read line by line, in \cref{wex:readwhile}.

\subsection{The clustering types: \texttt{Krylov}, \texttt{StepOutputs}, \texttt{PrevCarry}}

Before the operators come the records they thread, the ones the type-placement rule
keeps \emph{here} rather than in \emph{types} because only \emph{iteration} uses
them. \code{Krylov} is the ephemeral per-restart workspace---born at restart entry,
discarded at restart exit, threaded as a \emph{plain value} (not a monadic effect,
because its lifetime is strictly inside one cycle). Its fields differ by solver:

\begin{lstlisting}[style=l4style]
-- Ephemeral workspace; born at restart, discarded at restart exit. Plain value.
type Krylov_CG = {
  r  : Tensor[(S: ...)],  -- residual
  p  : Tensor[(S: ...)],  -- search direction
  z? : Tensor[(S: ...)],  -- preconditioned residual (iff a preconditioner is set)
  alpha : Scalar,         -- step length
  beta  : Scalar          -- direction-update coeff (the convergence-test proxy)
}
type Krylov_GMRES = {
  V  : [Tensor[(S: ...)]],  -- Arnoldi basis (array of basis columns)
  Z? : [Tensor[(S: ...)]],  -- preconditioned basis (iff OpParams.flexible -- FGMRES)
  H  : DenseMatrix,         -- upper-Hessenberg (small-dense, scalar-stratum)
  s, cs, sn : [Scalar],     -- LS RHS / Givens cosines / Givens sines (small-dense)
  beta : Scalar, j : Int    -- residual proxy; inner index within the restart cycle
}
\end{lstlisting}
The contrast is the whole story of CG versus GMRES: CG's workspace is a fixed
handful of vectors, so its memory is bounded; GMRES's \code{V} \emph{grows} with
each inner iteration, which is exactly why GMRES restarts. \code{StepOutputs} is
the demand-prunable per-step readout (a residual norm, an optional
least-squares residual, an optional breakdown token)---the ``extras'' of
\cref{ch:fold}, rendered as a record. And \code{PrevCarry} is the one-slot
recurrence carry that Form~B of the Krylov step threads through the loop instead of
storing in the state---the first-iteration-unrolling of \cref{ch:fold} made
concrete.

\subsection{\texttt{iterate\_while}: the def body itself}

Here is the combinator the whole book turns on, rendered as synthesized code. We
read it in full in \cref{wex:readwhile}; for now, see it whole.

\begin{lstlisting}[style=l4style]
-- # Arguments
--   a    : a          -- the value-threaded carry (immutable); init
--   cont : a -> Bool  -- loop predicate, PURE on the carry; fires before each step
--   step : a -> Solve { state: a, ...e }  -- per-step body; Solve effect on SimState
-- # Returns
--   Solve { final_state: a, trajectory: [{ ...e }] }
--     trajectory -- per-step extras in order; demand-pruned when only final_state read
iterate_while :: a -> (a -> Bool) -> (a -> Solve { state: a, ...e })
              -> Solve { final_state: a, trajectory: [{ ...e }] }
iterate_while a cont step =
  if cont a
    then do
      { state: a', ...e } <- step a
      { final_state, trajectory } <- iterate_while a' cont step
      pure { final_state, trajectory: [{ ...e }] ++ trajectory }
    else
      pure { final_state: a, trajectory: [] }

-- the no-extras sugar (a bounded loop): e = (), body non-monadic.
iterate_while_pure :: a -> (a -> Bool) -> (a -> a) -> a
iterate_while_pure a cont f =
  (iterate_while a cont (\x -> pure { state: f x })).final_state
\end{lstlisting}
This is the small-step rule of \cref{ch:fold} transcribed verbatim into code: test
\code{cont a}; if true, run one \code{step}, recurse on the new carry, and prepend
this step's extras to the trajectory; if false, stop with the carry as
\code{final\_state} and an empty trajectory. The recursion sits in tail position, so
it is a loop, not a stack. The \code{iterate\_while\_pure} below it is the sugar for
loops with nothing to report---the carry threads, the trajectory is always empty,
and you read back only the final value. \Cref{fig:whilestructure} diagrams the
structure of the def; \cref{wex:readwhile} reads it line by line.

\begin{figure}[t]
\centering
\begin{tikzpicture}[font=\footnotesize, node distance=5mm]
  % the predicate gate
  \node[diamond, draw=simBlue, fill=simBgBlue, aspect=1.7, inner sep=1pt,
    font=\scriptsize] (gate) at (0,0) {\code{cont a}?};
  % FALSE branch (else)
  \node[product, minimum width=30mm, align=left, right=20mm of gate] (base)
    {\scriptsize\textbf{else (base case)}\\\scriptsize \code{final\_state: a}\\\scriptsize \code{trajectory: []}};
  \draw[flow] (gate.east) -- node[above,font=\scriptsize,text=simRust]{no} (base.west);
  % TRUE branch (then) -- three lines stacked below
  \node[opbox, minimum width=46mm, align=left, below=10mm of gate] (s1)
    {\scriptsize 1.\ \code{\{state: a', ...e\} <- step a}\\\scriptsize\itshape one step: next carry + extras};
  \node[opbox, minimum width=46mm, align=left, below=4mm of s1] (s2)
    {\scriptsize 2.\ \code{... <- iterate\_while a' cont step}\\\scriptsize\itshape tail recurse on the new carry};
  \node[opbox, minimum width=46mm, align=left, below=4mm of s2] (s3)
    {\scriptsize 3.\ \code{pure \{..., trajectory: [e] ++ ...\}}\\\scriptsize\itshape prepend this step's extras};
  \draw[flow] (gate.south) -- node[left,font=\scriptsize,text=simGreen]{yes} (s1.north);
  \draw[flow] (s1) -- (s2);
  \draw[flow] (s2) -- (s3);
  % the carry feedback to the gate (recursion)
  \draw[backflow] (s2.east) .. controls +(2.4,0) and +(2.4,0) ..
    node[right,font=\scriptsize,text=simRust,align=left,pos=0.5]{recurse:\\carry \code{a'}} (gate.east);
\end{tikzpicture}
\caption[The structure of the \texttt{iterate\_while} def]{The synthesized
\code{iterate\_while} def as a structure diagram. The predicate \code{cont a} is
tested \emph{first} (the \texttt{while} convention of \cref{ch:fold}). On
\textcolor{simRust}{no}, the \texttt{else} arm returns the base case---the carry as
\code{final\_state}, an empty trajectory. On \textcolor{simGreen}{yes}, the
\texttt{then} arm runs three lines: take one step (producing the next carry \code{a'}
and this step's extras \code{e}), tail-recurse on \code{a'} (the dashed feedback
edge), and prepend \code{e} to the trajectory the recursion returns. Three lines and
a gate; \cref{wex:readwhile} reads each in turn.}
\label{fig:whilestructure}
\end{figure}

\subsection{The step kernels: \texttt{krylov\_step} and \texttt{chebyshev}}

Above the loop sit the kernels it folds. \code{krylov\_step} is the per-step body
of a Krylov solve, embedded in the \code{Solve} monad. It comes in two
\emph{forms}, and seeing them side by side is the clearest picture of the
first-iteration-unrolling rotation from \cref{ch:fold}. \emph{Form~A}---the
default---keeps the ``is this the first iteration?'' branch \emph{inside} the body
and folds with plain \code{iterate\_while}:

\begin{lstlisting}[style=l4style]
-- Form A (branch-in-body, default): folded by iterate_while.
krylov_step :: OpParams -> Krylov -> (SimState -> Solve { krylov: Krylov, outputs: StepOutputs })
krylov_step op K = \s -> do
  let w     = apply_linop op.T K.input_field   -- operator apply (no SimState read)
  let K_aux = optionally_apply_auxiliary op K w -- GMRES orthog / Chebyshev scalars / CG no-op
  let K'    = krylov_update K_aux op w          -- bundle update (axpy/dot/nrm2, pure on K)
  let outputs = derived_views K' op             -- residual_norm; ls_residual; breakdown
  modify (\s -> s { it = s.it + 1 })            -- the sole monadic effect
  pure { krylov: K', outputs }
\end{lstlisting}
\emph{Form~B}---opt-in---splits the body into a \code{first\_step} that
manufactures the initial recurrence value and a \emph{branch-free}
\code{steady\_step} that consumes it, folded by
\code{iterate\_while\_with\_prev}. The canonical instance is CG, and reading the
two CG steps next to each other shows what unrolling buys: the steady step's
\code{p'} line is unconditional, where Form~A's would carry an \code{if it==0}
test.

\begin{lstlisting}[style=l4style]
-- Form B, CG: the steady step is branch-free; beta_prev is the closure carry.
cg_steady_step :: LinOp[(S: ...)] -> Scalar -> Scalar -> CgState
              -> { state: CgState, residual_norm: Scalar }
cg_steady_step opA eps beta_prev s =
  let p'    = axpby 1.0 s.r (s.beta / beta_prev) s.p in  -- unconditional: p <- r + (beta/beta_prev) p
  let Ap    = apply_linop opA p' in
  let alpha = s.beta / dot Ap p' in
  let x'    = axpy alpha p' s.x in                       -- x <- x + alpha p
  let r'    = axpy (negate alpha) Ap s.r in              -- r <- r - alpha Ap
  let beta' = dot r' r' in
  let res'  = sqrt (abs beta') in
  { state: { x: x', r: r', p: p', beta: beta', it: s.it + 1, converged: res' < eps },
    residual_norm: res' }
\end{lstlisting}
Every line of that step is a \emph{data-algebra} call---\code{axpby}, \code{axpy},
\code{dot}---which is the cue for the next shelf. But notice first that the step's
shape is exactly the carry-threading fold of \cref{ch:fold}: a record in, a record
out, the termination quantity (\code{beta}, \code{converged}) folded into the
carry so the predicate can read it. \code{chebyshev}, the fourth operator, is the
same story at two nesting levels: an outer Richardson sweep and an inner
polynomial recurrence, both written as bounded \code{iterate\_while\_pure} folds
with step-count predicates. We do not reproduce its body here---it is long, and
\cref{ch:matrixfree} reads it as a smoother---but its shape is no surprise: a fold
inside a fold, each a setting of the one combinator.

\begin{notationbox}
The \emph{iteration} library has \emph{no} \code{\#extern} kernel of its own. The
loop combinators are pure calculus, and the opaque library boundaries the steps
eventually hit---the libCEED quadrature inside \code{apply\_linop}, the SLEPc
eigensolve loop---belong to the operators that \emph{own} those applies, and render
\code{\#extern} on the \emph{data-algebra} and \emph{coordination} shelves
instead. Here, \code{apply\_linop} is an abstract surface the step folds; its
implementation lives where the operator is built.
\end{notationbox}

\section{The \texttt{data-algebra} library: the tensor verbs}
\label{sec:dataalgebra}

The third shelf holds the verbs that \emph{combine} tensors---the right-hand branch
of the combinator vocabulary in \cref{ch:combinators}. Every line of the CG step we
just read was one of these. There are two combinators at the top
(\code{linear\_combination} and \code{inner\_product}), their named leaves and
consumers, an operator constructor, the assemble fold with its one \code{\#extern}
kernel, and a family of reduce verbs. We render the headline def and survey the
rest.

\subsection{\texttt{linear\_combination} and its arity leaves}

The most-used verb in the book is \code{linear\_combination}, the scalar-weighted
tensor sum $\sum_i a_i t_i$. It is a \code{foldl} over a term list, and the four
BLAS routines hang off it as \code{where}-local arity leaves---specializations at a
fixed term-list length, exactly the picture of \cref{fig:lincomb}.

\begin{lstlisting}[style=l4style]
-- The scalar-weighted-tensor-sum combinator: a pure fold over a (Scalar, Tensor)
-- term list, producing the sum of a_i * t_i. No Solve monad, no carry.
-- Laws (concatenation-homomorphism, multilinearity): see ch:combinators.
linear_combination :: [(Scalar, Tensor[(S: ...)])] -> Tensor[(S: ...)]
linear_combination pairs = foldl (\acc (a, t) -> acc + scal a t) (zeros) pairs
  where
    -- the four BLAS-1 arity leaves: this combinator at a fixed term-list length
    -- (specialization notes, not co-equal ops -- accelerated kernels stopped low).
    scal     a x          = linear_combination [(a, x)]
    axpy     a x y        = linear_combination [(a, x), (1, y)]
    axpby    a x b y      = linear_combination [(a, x), (b, y)]
    axpbypcz a x b y c z  = linear_combination [(a, x), (b, y), (c, z)]
\end{lstlisting}
Read the four leaves top to bottom: \code{scal} is a one-term list, \code{axpy} a
two-term list with the second coefficient pinned to $1$, \code{axpby} the general
two-term list, \code{axpbypcz} the three-term list. They differ only in arity, and
the concatenation law makes the chain exact---\code{axpbypcz}'s list is
\code{axpby}'s list concatenated with \code{scal}'s. \Cref{fig:lincombleaves}
redraws this as the trunk-and-leaves the synthesized code makes literal: one def, a
\code{where}-block of four one-liners.

\begin{figure}[t]
\centering
\begin{tikzpicture}[font=\footnotesize, node distance=5mm]
  % trunk
  \node[stage, minimum width=66mm, align=center] (trunk) at (0,0)
    {\code{linear\_combination} \;\itshape\quad sum of $a_i\,t_i$, term list, any length};
  \def\ly{-1.7}
  \node[opbox, minimum width=24mm, align=center] (l1) at (-4.8,\ly)
    {\code{scal}\\\scriptsize arity 1\\\scriptsize \code{[(a,x)]}};
  \node[opbox, minimum width=24mm, align=center] (l2) at (-1.6,\ly)
    {\code{axpy}\\\scriptsize arity 2\\\scriptsize \code{[(a,x),(1,y)]}};
  \node[opbox, minimum width=24mm, align=center] (l3) at (1.6,\ly)
    {\code{axpby}\\\scriptsize arity 2\\\scriptsize \code{[(a,x),(b,y)]}};
  \node[opbox, minimum width=24mm, align=center] (l4) at (4.8,\ly)
    {\code{axpbypcz}\\\scriptsize arity 3\\\scriptsize \code{[..,(c,z)]}};
  \foreach \l in {l1,l2,l3,l4}{ \draw[refedge] (trunk.south) -- (\l.north); }
  \node[font=\scriptsize\itshape, text=simGray, align=center] at (0,-2.9)
    {rendered as a \code{where}-block of four one-liners beneath the one def:\\
     the combinator is the entry; the leaves are settings of it};
\end{tikzpicture}
\caption[\texttt{linear\_combination} with its \texttt{where}-local leaves]{The
synthesized \code{linear\_combination} and its four BLAS arity leaves, as the code
renders them: one def for the combinator (trunk), a \code{where}-block of four
one-liners beneath it (leaves). Each leaf is the combinator at a fixed term-list
length---\code{scal} at one term, \code{axpy}/\code{axpby} at two, \code{axpbypcz}
at three. This is the combinator-primary model of \cref{ch:combinators} made
literal in the library: the fused kernels are ``stopped low'' as notes, and the
combinator rises to be the entry. Compare \cref{fig:lincomb}.}
\label{fig:lincombleaves}
\end{figure}

\subsection{\texttt{inner\_product}, \texttt{dot}, and \texttt{nrm2}}

The scalar-producing sibling is \code{inner\_product}, $\langle x, y\rangle_M =
x^{\!\top} M y$, a reduction over the shape group. \code{dot} is its
$M = I$ \emph{leaf} (a specialization, hung below), and \code{nrm2} is a
\emph{consumer}---$\sqrt{\,\cdot\,}\circ\code{dot}$ at the diagonal, built
\emph{above} \code{dot}, not beneath it. The synthesized code keeps the distinction
visible:

\begin{lstlisting}[style=l4style]
inner_product :: Tensor[(S: ...)] -> Tensor[(S: ...)] -> Scalar
inner_product x y = reduce (+) zero (zipWith kernel x y)
  where kernel xi yi = conj_if_complex xi * yi   -- Hermitian: arg-1 conjugated

dot  x y = inner_product x y                       -- the M = I leaf
nrm2 :: Tensor[(S: ...)] -> Scalar
nrm2 x = sqrt (abs (inner_product x x))            -- a CONSUMER of inner_product
\end{lstlisting}
The \code{abs} inside \code{nrm2} is not decoration---it is the defensive
non-negativity guard that keeps a round-off-negative radicand from reaching the
square root, a load-bearing numerical detail preserved as an explicit part of the
scalar map (\cref{ch:bigidea}). The CG step's \code{dot Ap p'} and the residual's
\code{sqrt (abs beta')} we read in \cref{sec:iteration} are exactly these verbs.

\subsection{The operator constructor and the assemble fold}

Two larger items sit on this shelf. \code{mk\_matrix\_free\_operator} builds an
operator \emph{value}---a closure whose \code{apply} is the five-stage contraction
chain $G^{\!\top}\!\cdot B^{\!\top}\!\cdot D \cdot B \cdot G$ of \cref{ch:matrixfree},
rendered inline because that chain \emph{is} the implementation:

\begin{lstlisting}[style=l4style]
-- A = G^T . B_D^T . D . B_D . G  (the un-materialized contraction chain; ch:matrixfree)
mk_matrix_free_operator space term geom = mkOp (\v -> apply_chain v)
  where apply_chain v =
    element_restrict_T (basis_apply_T term
      (quad_point_contract geom (basis_apply term (element_restrict space v))))
\end{lstlisting}
And \code{fe\_assemble} folds the weak-form-term list into a global operator,
$K = \sum_t \code{assemble\_term}(\text{space}, t)$---a \code{foldr} over a
commutative-monoid sum. Its per-term leaf is the \emph{first} \code{\#extern} we
meet: \code{assemble\_term} is the libCEED-owned quadrature kernel, rendered as a
callout in place of a body.

\begin{lstlisting}[style=l4style]
fe_assemble space terms = foldr (\t acc -> assemble_term space t + acc) zero terms

-- The opaque per-term leaf: ONE weak-form term to its global-dof contribution.
-- libCEED-owned (the kernel-API boundary). Rendered #extern after its signature.
assemble_term :: FiniteElementSpace[N] -> WeakFormTerm -> LinOp[(N: ...)]
#extern assemble_term
\end{lstlisting}

\subsection{The reduce verbs}

The bottom of the shelf holds the reductions of \cref{ch:reductions}---the verbs
that turn a solved field into the physical product. There are five firm ones, each
a \code{map}-then-collect over a family with no inter-element state:
\code{gram\_reduce} (the symmetric Gram matrix behind capacitance and inductance),
\code{domain\_energy\_reduce} (the per-domain energy table), \code{sparameter\_reduce}
(the driven scattering matrix), \code{eigenfreq\_qfactor\_reduce} (the per-mode
frequency and quality factor), and \code{waveguide\_mode\_reduce} (the boundary-mode
field table). They share a shape but deliberately not a definition---a scalar table,
a matrix, and a field-carrying table are different products, and \cref{ch:reductions}
explains why forcing them together would be the wrong design. One example shows the
shape:

\begin{lstlisting}[style=l4style]
-- per-mode (f, Q) table over the converged eigenpairs: map-then-collect, no state.
eigenfreq_qfactor_reduce ptype kappa eigs =
  [ let omega = untransform ptype lambda            -- eigenvalue -> angular frequency
        f     = re omega                            -- eigenfrequency f = Re omega
        k     = kappa mode                          -- loss rate
        q     = if k == 0 then infinity else f / abs k  -- quality factor Q = omega / kappa
    in (f, q)
  | (lambda, _E) <- eigs ]                          -- map over the eigenpair family
\end{lstlisting}
The \code{if k == 0 then infinity} is the lossless-cavity guard---a cavity with no
loss has infinite $Q$---the same kind of explicit boundary case as \code{nrm2}'s
\code{abs}. Read past the physics and the shape is a \code{map} with a guard: the
combinator-primary model again, this time on the reduce side.

\section{The \texttt{coordination} library: the outer caps}
\label{sec:coordination}

The fourth shelf holds the \code{Solve}-monadic outer caps---the combinators that
drive the iteration kernels to convergence and run them over families. This is the
left-hand branch of \cref{ch:combinators}'s vocabulary and the home of the monad of
\cref{ch:monad}. Its headline def, \code{ksp\_solve}, is the second we read line by
line, in \cref{wex:kspsolve}.

\subsection{The \texttt{Solve} monad and the termination sums}

The coordination surface threads \code{SimState} through a state monad and
classifies termination into sum types. The monad is the thin one of \cref{ch:monad}:

\begin{lstlisting}[style=l4style]
-- The outer-driver state monad: the effect domain is exactly SimState.
type Solve a = StateT SimState Identity a

-- the discharge that turns a monadic cap into a pure function:
execState :: Solve a -> SimState -> SimState   -- run the action, project the state
\end{lstlisting}
And termination is a \emph{value}, a sum type the cap classifies once per cycle.
Krylov solves use a three-arm \code{Outcome}; the eigensolver uses a richer four-arm
\code{EigOutcome} that adds a first-class ``partially converged'' arm with no Krylov
analog:

\begin{lstlisting}[style=l4style]
data Outcome    = Continue | Done Bool                    -- the Krylov 3-arm sum
data EigStatus  = Converged | PartialConverged Int | MaxIterReached | LinearSolveFailed
data EigOutcome = Continue | Done EigStatus               -- the eigen 4-arm extension
\end{lstlisting}
The eigensolver's persistent state, \code{EigState}, is a single-library record
(only \code{eigsolve} names it), so the type-placement rule keeps it here, not in
\emph{types}. It does not collapse to \code{SimState} because it carries a whole
\emph{family} of converged eigenpairs, not one iterate.

\subsection{\texttt{ksp\_solve}: the Krylov cap}

Here is the cap, the def \cref{wex:kspsolve} reads. It seeds a \code{SimState},
runs the outer driver over it, and projects the terminal state with
\code{execState}---turning a monadic computation back into a pure
$\code{OpParams}\to\code{Inputs}\to\code{SimState}$ function.

\begin{lstlisting}[style=l4style]
-- # Returns SimState: terminal .x ~ A^-1 b, plus .it / .converged / .final_res
ksp_solve :: OpParams -> Inputs -> SimState
ksp_solve op inp = execState (solve_loop op inp) (initial_state inp)
  where
    -- outer driver: tail-recurse the per-cycle body until the Outcome says stop.
    solve_loop op inp = do
      o <- restart_cycle op inp
      unless (done o) (solve_loop op inp)

    -- per-cycle body: fresh Krylov, inner kernel-fold, fold the correction, classify once.
    restart_cycle op inp = do
      s <- get
      let k0        = fresh_krylov op inp s                 -- ephemeral bundle (plain value)
          (kn, _os) = iterate_while (krylov_step op k0) cont -- inner fold of the kernel
      modify (\s -> s { x = s.x `plus` (kn.basis `applyBasis` kn.y) })  -- the single per-cycle x write
      pure (classify kn op s)                               -- Done True / Done False / Continue
\end{lstlisting}
The inner \code{krylov\_step} and \code{iterate\_while} are not re-rendered
here---they live on the \emph{iteration} shelf, and the cap simply \emph{folds}
them. That deep link is the topological order paying off: \code{coordination} sits
above \code{iteration}, so it calls down to it.

\subsection{\texttt{eigsolve} and the second \texttt{\#extern}}

\code{eigsolve} is the cap that does \emph{not} own its loop. Where \code{ksp\_solve}
writes the tail-recursion itself, the eigen-iteration runs \emph{inside} a library
(SLEPc's \code{EPSSolve}), so the cap names the iteration by role and renders it
\code{\#extern}---the second kernel boundary in the library:

\begin{lstlisting}[style=l4style]
eigsolve op inp = execState (solve_loop op inp) (initial_eig_state inp)
  where solve_loop op inp = do
    o <- eigen_iterate op inp     -- OPAQUE library fold (no Palace loop)
    classify_into_state o

-- The opaque eigen-iteration: the entire SLEPc EPSSolve loop. The kernel-API boundary.
eigen_iterate :: OpParams -> Inputs -> Solve EigOutcome
#extern eigen_iterate
\end{lstlisting}
This is the deliberate non-loop of \cref{ch:combinators}: the fourth solve corner,
the one outside the map/fold grid because we write no loop at all. The
\code{\#extern} marks exactly where our vocabulary stops and the library's begins.
(\Cref{ch:krylov} separately reconstructs what that loop does internally, in our own
\code{lanczos\_step}/\code{krylov\_step} vocabulary---a glass-box version that
\emph{realizes} this API for the curious reader, without the cap depending on it.)

\subsection{\texttt{solve\_family}, \texttt{frequency\_sweep}, \texttt{fold\_solve}}

The top of the shelf holds the three family combinators of \cref{ch:combinators},
and each is a near one-liner over the caps below it. \Cref{fig:caps} lines all five
coordination caps up as the map/fold/black-box family they form.

\begin{lstlisting}[style=l4style]
-- fixed-operator map: capture op once, map ksp_solve over the RHS family.
solve_family op rhss = map (\inp -> ksp_solve op inp) rhss

-- operator-VARYING map: rebuild A(omega) INSIDE the map, then solve.
frequency_sweep fam omegas =
  map (\w -> ksp_solve (assemble_frequency_operator fam w) (rhs_at fam w)) omegas

-- state-threaded fold: thread the field-state through a schedule by foldl.
fold_solve op s0 schedule = foldl (\s t -> time_step_op op s t) s0 schedule
time_step_op :: OpParams -> TimeState -> Time -> TimeState
#extern time_step_op                              -- the opaque per-step integrator
\end{lstlisting}
The three differ by exactly the axes \cref{ch:combinators} named:
\code{solve\_family} hoists the operator out of the map; \code{frequency\_sweep}
rebuilds it inside; \code{fold\_solve} drops the independence and threads a carry,
folding a third \code{\#extern} per-step integrator. Three lines, three corners.

\begin{figure}[t]
\centering
\begin{tikzpicture}[font=\footnotesize]
  % the family header
  \node[stage, minimum width=64mm, align=center] (hdr) at (0,2.0)
    {the \emph{coordination} caps: \code{map} / \code{fold} / black-box};
  % map row
  \node[opbox, minimum width=30mm, fill=simBgBlue, draw=simBlue, align=center] (sf)
    at (-3.6,0.3) {\code{solve\_family}\\\scriptsize fixed-op \textbf{map}};
  \node[opbox, minimum width=30mm, fill=simBgBlue, draw=simBlue, align=center] (fs)
    at (-0.0,0.3) {\code{frequency\_sweep}\\\scriptsize op-varying \textbf{map}};
  \node[opbox, minimum width=30mm, fill=simBgGreen, draw=simGreen, align=center] (fo)
    at (3.6,0.3) {\code{fold\_solve}\\\scriptsize state-threaded \textbf{fold}};
  % bottom row: the two caps the families call / the black box
  \node[opbox, minimum width=30mm, align=center] (ksp) at (-1.8,-1.6)
    {\code{ksp\_solve}\\\scriptsize the Krylov cap (our loop)};
  \node[extern, minimum width=30mm, align=center] (eig) at (1.8,-1.6)
    {\code{eigsolve}\\\scriptsize black-box (\#extern loop)};
  \draw[depedge] (hdr) -- (sf); \draw[depedge] (hdr) -- (fs); \draw[depedge] (hdr) -- (fo);
  \draw[refedge] (sf.south) -- (ksp.north);
  \draw[refedge] (fs.south) -- (ksp.north);
  \node[font=\scriptsize\itshape, text=simViolet] at (1.8,-2.6) {(no map/fold: a single call)};
\end{tikzpicture}
\caption[The coordination caps as the map/fold/black-box family]{The five
coordination caps as the family of \cref{ch:combinators}. \code{solve\_family} and
\code{frequency\_sweep} are \code{map}s---independent solves, the first hoisting its
operator, the second rebuilding it per member---and both call the Krylov cap
\code{ksp\_solve} (our loop) per element. \code{fold\_solve} is the
\code{fold}---a threaded carry, intrinsically sequential. \code{eigsolve} is the
black-box: a single \code{\#extern} library call, outside the map/fold grid because
no loop is ours to write. One header, three structural shapes.}
\label{fig:caps}
\end{figure}

\section{The same shapes, again and again}
\label{sec:samesahapes}

Step back from the four shelves and look at what we actually rendered. The
\emph{iteration} library was a fold and two kernels that are folds. The
\emph{data-algebra} library was two folds (\code{linear\_combination},
\code{inner\_product}) with named leaves, an operator that is a contraction chain,
and five reductions that are \code{map}s-then-collect. The \emph{coordination}
library was a state monad, three caps that are \code{map}s or a \code{fold}, and
two \code{\#extern} boundaries. And the \emph{types} library was three records,
distinguished only by lifetime. \Cref{fig:samesahapes} tallies the recurring
shapes across the libraries; the count is strikingly small.

\begin{figure}[t]
\centering
\setlength{\tabcolsep}{5pt}
\renewcommand{\arraystretch}{1.25}
\begin{tabular}{@{}lcccc@{}}
\toprule
& \textbf{types} & \textbf{iteration} & \textbf{data-algebra} & \textbf{coordination} \\
\midrule
a \emph{fold}        & --- & \code{iterate\_while} & \code{linear\_combination} & \code{fold\_solve} \\
                     &     & \code{krylov\_step}   & \code{inner\_product}      &                    \\
a \emph{map}         & --- & ---                   & the reduce verbs           & \code{solve\_family} \\
                     &     &                       &                            & \code{frequency\_sweep} \\
a record by lifetime & all 3 & \code{Krylov} etc.  & \code{DofSet} etc.         & \code{EigState} etc. \\
an \code{\#extern}   & --- & ---                   & \code{assemble\_term}      & \code{eigen\_iterate} \\
                     &     &                       &                            & \code{time\_step\_op} \\
\bottomrule
\end{tabular}
\caption[The recurring shapes, tallied across the libraries]{The four libraries,
tallied by the shape of what they hold. Every operator is a \emph{fold}, a
\emph{map}, or one of three \code{\#extern} boundaries; every type is a record
distinguished by its lifetime. The table is mostly the same four rows repeating
across the columns---which is the chapter's whole point. There is no fifth kind of
thing. The library is small because the calculus is small: a few combinators, a few
records threaded by lifetime, and a marked boundary where an opaque kernel sits.}
\label{fig:samesahapes}
\end{figure}

\begin{keyidea}
Across all four libraries you keep meeting the \emph{same handful of shapes}: a
fold (\code{iterate\_while}, \code{linear\_combination}, \code{fold\_solve}), a map
(\code{solve\_family}, the reduce verbs), a record distinguished only by its
lifetime (\code{IoData}/\code{OpParams}/\code{SimState} and their single-library
cousins), and a marked \code{\#extern} boundary where an opaque kernel lives
(\code{assemble\_term}, \code{eigen\_iterate}, \code{time\_step\_op}). There is no
fifth kind of thing. \textbf{The library is small because the calculus is small}:
the entire synthesized implementation of a hundred-thousand-line simulator is a few
combinators, a few records, and three kernel callouts---the vocabulary of
\cref{part:framework}, rendered as code.
\end{keyidea}

This is the dividend the whole book has been compounding. \Cref{part:framework}
argued that the simulator is built from a small calculus; \cref{part:machinery} and
\cref{part:modalities} spent the bulk of the text showing that the argument holds
component by component. Here, with the spec rendered as a library, we can simply
\emph{count}: four shelves, and on them a fold, a map, a record, and an
\code{\#extern}, repeated. The next chapter, \cref{ch:composing}, opens the fifth
shelf---\emph{drivers}---and watches a single analysis composed end to end out of
the items we have just toured. Nothing new appears there either; a driver is these
four libraries, wired together.

\section{Worked examples}

\begin{workedexample}{Reading the \texttt{iterate\_while} def body line by line}{readwhile}
We read the synthesized \code{iterate\_while} of \cref{sec:iteration} one line at a
time, connecting each to a concept from \cref{ch:fold}. The body is:
\begin{lstlisting}[style=l4style]
iterate_while a cont step =
  if cont a                                              -- (1)
    then do
      { state: a', ...e } <- step a                      -- (2)
      { final_state, trajectory } <- iterate_while a' cont step  -- (3)
      pure { final_state, trajectory: [{ ...e }] ++ trajectory }  -- (4)
    else
      pure { final_state: a, trajectory: [] }            -- (5)
\end{lstlisting}

\emph{Line (1): the predicate fires first.} \code{if cont a} tests the predicate
\emph{before} any step runs---this is the \texttt{while} convention of
\cref{ch:fold}, Rule~1. If \code{cont a} is already false on the initial carry, the
body never executes and we fall straight to line~(5). A solver handed an
already-converged input does zero work, exactly as promised.

\emph{Line (2): one step, producing the next carry and the extras.}
\code{\{ state: a', ...e \} <- step a} runs the step once and \emph{destructures}
its output record: the \code{state} field becomes the next carry \code{a'}, and the
remaining fields---the open \code{...e}---are this step's extras. The \code{<-} is
the monadic bind: \code{step} runs in the \code{Solve} monad, so this line also
threads the \code{SimState} effect (the counter increment) under the hood, while
the carry \code{a'} stays a pure value. This is the carry/extras split of
\cref{ch:fold}: the carry threads forward, the extras peel off.

\emph{Line (3): tail-recurse on the new carry.}
\code{iterate\_while a' cont step} calls the combinator again with the
\emph{new} carry \code{a'} and the \emph{same} predicate and step. Because this call
is the last thing the \texttt{then} arm does before line~(4) assembles the result,
it is in tail position---the implementation realizes it as a loop, not a growing
stack. The predicate will be re-tested on \code{a'} at the top of this recursive
call; we may \emph{not} hoist that test out (\cref{ch:fold}'s non-law), or the loop
never stops.

\emph{Line (4): prepend this step's extras to the trajectory.}
\code{[\{ ...e \}] ++ trajectory} puts this step's extras at the \emph{front} of
whatever trajectory the recursion returned. Unrolled, the extras line up in
iteration order: step~1's first, step~2's next, and so on---the trajectory of
\cref{fig:unroll}. And because the trajectory is a \emph{value} the consumer may or
may not read, the demand-driven pruning of \cref{ch:fold} applies: if nothing
downstream reads \code{trajectory}, the \code{e}-producing work in line~(2) is
eliminated at every step, and what remains is the bare carry spine of
\code{iterate\_while\_pure}.

\emph{Line (5): the base case.} \code{\{ final\_state: a, trajectory: [] \}}
returns the current carry as the answer with an empty trajectory. This is the
\texttt{else} arm, reached the first time \code{cont a} is false. It is the formal
content of \cref{ch:fold}'s ``base case'' law: for \emph{any} step, an initial
carry that fails the predicate yields itself and nothing else.

Five lines: a gate, three lines of recursive work, one base case. Every iterative
algorithm in the simulator---CG, GMRES, time-stepping, the adaptive
loop---specializes \emph{this} body by supplying its own carry type, predicate, and
step. The loop is the same five lines every time.
\end{workedexample}

\begin{workedexample}{Reading the \texttt{ksp\_solve} cap}{kspsolve}
We read the \code{ksp\_solve} of \cref{sec:coordination}, connecting it to the monad
of \cref{ch:monad} and the Krylov machinery of \cref{ch:krylov}. The cap is:
\begin{lstlisting}[style=l4style]
ksp_solve op inp = execState (solve_loop op inp) (initial_state inp)  -- (1)
  where
    solve_loop op inp = do
      o <- restart_cycle op inp                            -- (2)
      unless (done o) (solve_loop op inp)                  -- (3)
    restart_cycle op inp = do
      s <- get                                             -- (4)
      let k0        = fresh_krylov op inp s                -- (5)
          (kn, _os) = iterate_while (krylov_step op k0) cont  -- (6)
      modify (\s -> s { x = s.x `plus` (kn.basis `applyBasis` kn.y) })  -- (7)
      pure (classify kn op s)                              -- (8)
\end{lstlisting}

\emph{Line (1): the \code{execState} discharge.} The cap's type is
$\code{OpParams}\to\code{Inputs}\to\code{SimState}$---a \emph{pure} function---yet
its body is monadic. \code{execState} bridges the two: it runs the
\code{Solve}-action \code{solve\_loop op inp} starting from a seeded state
\code{initial\_state inp}, threads the \code{SimState} through every \code{modify}
inside, and projects out the final state, dropping the action's unit return value.
This is exactly the monad-to-pure-function discharge of \cref{ch:monad}: the thin
\code{StateT} wrapper exists only inside the cap, and \code{execState} peels it
away at the boundary so the outside world sees a plain function.

\emph{Lines (2)--(3): the outer restart loop.} \code{solve\_loop} runs one
\code{restart\_cycle}, binds its \code{Outcome} to \code{o}, and---via
\code{unless (done o)}---recurses only if the outcome is not \code{Done}. This is
the GMRES \emph{restart} structure of \cref{ch:krylov}: an \emph{outer} loop around
the inner Krylov fold, deliberately not flattened (the two stop on different
conditions, the non-fusion non-law of \cref{ch:fold}). For CG and Chebyshev, which
do not restart, this outer loop runs exactly once.

\emph{Line (4): read the state.} \code{s <- get} reads the current \code{SimState}
out of the monad---the iterate, the counter, the residuals as they stand at the
start of this cycle.

\emph{Line (5): a fresh ephemeral workspace.} \code{fresh\_krylov op inp s} builds
the \code{Krylov} bundle of \cref{sec:iteration}---born here, at cycle entry, as a
\emph{plain value}, not a monadic effect. Its lifetime is exactly this one cycle;
it will be discarded when the cycle ends. This is the ephemeral stratum of
\cref{ch:strata} in action.

\emph{Line (6): the inner Krylov fold.} \code{iterate\_while (krylov\_step op k0) cont}
is the heart---the very combinator we read in \cref{wex:readwhile}, now folding the
\code{krylov\_step} kernel of \cref{sec:iteration} over the working state until the
convergence predicate \code{cont} fails. Every CG/GMRES iteration of
\cref{ch:krylov} happens inside this one call. The result \code{kn} is the final
Krylov bundle; the discarded \code{\_os} would be its trajectory, pruned here
because the cap does not read it.

\emph{Line (7): the single per-cycle \code{x} write.} \code{modify (\textbackslash s -> s \{ x = ... \})}
is the cap's \emph{only} write to the iterate. After the inner fold finishes, the
accumulated correction \code{kn.basis `applyBasis` kn.y} (the back-substituted
solution increment) is added to \code{x} once. Not per inner iteration---\emph{once
per restart cycle}. This is the deliberate write discipline of \cref{ch:monad}: the
externally-visible iterate changes at cycle boundaries, while the inner fold works
in the ephemeral \code{Krylov} space and never touches \code{SimState.x}.

\emph{Line (8): classify the outcome once.} \code{classify kn op s} inspects the
final bundle (its residual proxy, its inner-iteration count) and the state, and
returns \code{Done True} (converged), \code{Done False} (hit \code{max\_it}), or
\code{Continue} (restart). The classification happens \emph{once}, at the cycle
boundary---the \code{Outcome} sum of \cref{sec:coordination}, the
classify-once law of \cref{ch:monad}.

Read whole: \code{ksp\_solve} is an outer restart loop (lines 2--3) around an inner
Krylov fold (line 6), with one state read (4), one ephemeral workspace (5), one
iterate write per cycle (7), and one termination classification (8), all discharged
to a pure function by \code{execState} (1). Every Krylov solver in the book---CG,
GMRES, FGMRES---is this same cap; what changes is only the \code{krylov\_step}
kernel it folds.
\end{workedexample}

\summarysection
The synthesized library is five modules in topological order; this chapter toured
four of them. The \emph{types} library holds three cross-cutting records
distinguished by lifetime: \code{IoData} (parsed once, selects the driver via
\code{problem.type}), \code{OpParams} (fixed at build time, variant selectors
folded into surfaces), and \code{SimState} (the run-time state, uniform across all
Krylov solvers). The \emph{iteration} library holds the \code{iterate\_while}
combinator of \cref{ch:fold}---rendered here as five lines of code---plus the
\code{krylov\_step} kernel (Form~A branch-in-body, Form~B unrolled) and the
\code{chebyshev} smoother, with \code{Krylov}/\code{StepOutputs}/\code{PrevCarry} as
single-library types. The \emph{data-algebra} library holds
\code{linear\_combination} with its \code{where}-local BLAS arity leaves,
\code{inner\_product}/\code{dot}/\code{nrm2}, the matrix-free operator constructor,
the \code{fe\_assemble} fold with its libCEED \code{\#extern}, and the five reduce
verbs. The \emph{coordination} library holds the \code{Solve} monad, the
\code{ksp\_solve} cap (read line by line), the \code{eigsolve} black-box with its
SLEPc \code{\#extern}, and the three family combinators
\code{solve\_family}/\code{frequency\_sweep}/\code{fold\_solve}. Across all four,
the same handful of shapes recurs---a fold, a map, a record threaded by lifetime, a
marked \code{\#extern} boundary---because the library is small for the same reason
the calculus is small. The fifth library, \emph{drivers}, composes these four and
is \cref{ch:composing}.

\furtherreading
The implementation idioms rendered here---the state monad as a thin wrapper
discharged by \code{execState}, \code{where}-local helper clusters, the records
threaded by lifetime---are standard functional-programming practice; Peyton Jones's
account of Haskell's design~\citep{peytonjones2003} is the canonical reference for
the monadic and record machinery. The Krylov kernels the \emph{iteration} and
\emph{coordination} libraries fold (CG, GMRES, and the Arnoldi/Lanczos relatives
behind \code{eigsolve}) are developed in full by Saad~\citep{saad2003}. The libCEED
quadrature kernel behind \code{fe\_assemble}'s \code{\#extern} boundary---the
matrix-free element-local contraction that \code{mk\_matrix\_free\_operator} renders
inline---is documented by the CEED project~\citep{ceed2021}.

\exercisessection
\begin{enumerate}[nosep]
  \item For each of the three \emph{types} records, name the stratum it belongs to
  (\cref{ch:strata}) and one field whose value is fixed for the whole run. Which of
  the three is the \emph{only} one written during the solve?
  \item \code{SimState} has the same five fields for CG and for GMRES, yet the two
  solvers clearly differ. Where does the difference live? (Hint: which library, and
  which record, holds the search direction for CG versus the growing basis for
  GMRES?)
  \item In the \code{iterate\_while} body of \cref{wex:readwhile}, suppose a
  consumer reads only \code{final\_state} and never the \code{trajectory}. Which
  line's work is pruned, and what combinator does the def collapse to? (Cite the
  relevant law from \cref{ch:fold}.)
  \item The four BLAS leaves of \code{linear\_combination} are written as a
  \code{where}-block of one-liners. Write \code{axpbypcz} in terms of \code{axpby}
  and \code{scal} using the concatenation law, and explain why this makes them ``one
  combinator, four leaves'' rather than four separate operators (\cref{ch:combinators}).
  \item \code{nrm2} calls \code{inner\_product} but is \emph{not} a leaf of it,
  while \code{dot} \emph{is}. Explain the difference using the
  specialization-versus-consumer distinction of \cref{ch:combinators}, and say what
  the \code{abs} inside \code{nrm2} guards against (\cref{ch:bigidea}).
  \item Three operators across the four libraries render \code{\#extern}:
  \code{assemble\_term} (\emph{data-algebra}), and \code{eigen\_iterate} /
  \code{time\_step\_op} (\emph{coordination}). For each, name the opaque library
  that owns it and say in one sentence why the synthesized code marks a boundary
  there rather than rendering a body.
  \item Reading \code{ksp\_solve} (\cref{wex:kspsolve}): the cap writes
  \code{SimState.x} exactly once per restart cycle (line~7), not once per inner
  iteration. Where do the \emph{inner} iterations do their work instead, and why is
  it correct for the externally-visible iterate to change only at cycle boundaries?
  Relate your answer to the strata of \cref{ch:strata}.
  \item \code{solve\_family}, \code{frequency\_sweep}, and \code{fold\_solve} are
  each near one-liners. Which two are \code{map}s and which is a \code{fold}? For the
  two maps, what single difference distinguishes them (\cref{ch:combinators})? Which
  of the three forbids reordering its elements, and why?
\end{enumerate}
